\documentclass[UTF8,12pt]{ctexart}
\usepackage[a4paper,margin=24mm]{geometry}
\usepackage{amsmath,amssymb,bm,graphicx,adjustbox}
\newcommand{\fitmath}[1]{\begin{adjustbox}{max width=\linewidth}$\displaystyle #1$\end{adjustbox}}
\setlength{\parindent}{0pt}
\setlength{\parskip}{.7em}
\sloppy
\allowdisplaybreaks
\begin{document}
\section*{1994 年考研数学三 · 真题}
\subsection*{1994 年全国硕士研究生招生考试试题（试卷 IV）}

\subsection*{一、填空题（本题共 5 小题，每小题 3 分，满分 15 分）}

（1）\(\displaystyle \int_{-2}^2\dfrac{\displaystyle x+|x|}{\displaystyle 2+x^2}\,dx=\)\_\_\_\_\_\_。


（2）已知 \(\displaystyle f^{\prime}(x_0)=-1\)，则 \(\displaystyle \lim_{x\to0}\dfrac{\displaystyle x}{\displaystyle f(x_0-2x)-f(x_0-x)}=\)\_\_\_\_\_\_。


（3）设方程 \(\displaystyle e^{xy}+y^2=\cos x\) 确定 \(\displaystyle y\) 为 \(\displaystyle x\) 的函数，则 \(\displaystyle \dfrac{\displaystyle dy}{\displaystyle dx}=\)\_\_\_\_\_\_。


（4）设 \(\displaystyle A=\begin{pmatrix}0&a_1&0&\cdots&0\\0&0&a_2&\cdots&0\\\vdots&\vdots&\vdots&&\vdots\\0&0&0&\cdots&a_{n-1}\\a_n&0&0&\cdots&0\end{pmatrix}\)，其中 \(\displaystyle a_i\ne0,\allowbreak  i=1,\allowbreak 2,\allowbreak \ldots,\allowbreak n\)，则 \(\displaystyle A^{-1}=\)\_\_\_\_\_\_。


（5）设随机变量 \(\displaystyle X\) 的概率密度为 \(\displaystyle f(x)=\begin{cases}2x,\allowbreak &0<x<1,\allowbreak \\0,\allowbreak &\text{其他},\allowbreak \end{cases}\) 以 \(\displaystyle Y\) 表示对 \(\displaystyle X\) 的三次独立重复观察中事件 \(\displaystyle \{X\leq\dfrac{\displaystyle 1}{\displaystyle 2}\}\) 出现的次数，则 \(\displaystyle P\{Y=2\}=\)\_\_\_\_\_\_。


\subsection*{二、选择题（本题共 5 小题，每小题 3 分，满分 15 分）}

（1）曲线 \(\displaystyle y=e^{1/x^2}\arctan\dfrac{\displaystyle x^2+x-1}{\displaystyle (x+1)(x-2)}\) 的渐近线有（　）。


（A）1 条。 （B）2 条。 （C）3 条。 （D）4 条。


（2）设常数 \(\displaystyle \lambda>0\)，而级数 \(\displaystyle \sum_{n=1}^{\infty}a_n^2\) 收敛，则级数 \(\displaystyle \sum_{n=1}^{\infty}(-1)^n\dfrac{\displaystyle |a_n|}{\displaystyle \sqrt{n^2+\lambda}}\)（　）。


（A）发散。 （B）条件收敛。 （C）绝对收敛。 （D）敛散性与 \(\displaystyle \lambda\) 有关。


（3）设 \(\displaystyle A\) 是 \(\displaystyle m\times n\) 矩阵，\(\displaystyle C\) 是 \(\displaystyle n\) 阶可逆矩阵，矩阵 \(\displaystyle A\) 的秩为 \(\displaystyle r\)，矩阵 \(\displaystyle B=AC\) 的秩为 \(\displaystyle r_1\)，则（　）。


（A）\(\displaystyle r>r_1\)。 （B）\(\displaystyle r<r_1\)。 （C）\(\displaystyle r=r_1\)。 （D）\(\displaystyle r\) 与 \(\displaystyle r_1\) 的关系依 \(\displaystyle C\) 而定。


（4）设 \(\displaystyle 0<P(A)<1,\allowbreak 0<P(B)<1\)，\(\displaystyle P(A\mid B)+P(\bar A\mid\bar B)=1\)，则事件 \(\displaystyle A\) 和 \(\displaystyle B\)（　）。


（A）互不相容。 （B）互相对立。 （C）不独立。 （D）独立。


（5）设 \(\displaystyle X_1,\allowbreak X_2,\allowbreak \ldots,\allowbreak X_n\) 是来自正态总体 \(\displaystyle N(\mu,\allowbreak \sigma^2)\) 的简单随机样本，\(\displaystyle \bar X\) 是样本均值，记 \(\displaystyle S_1^2=\dfrac{\displaystyle 1}{\displaystyle n-1}\displaystyle \sum_{i=1}^n(X_i-\bar X)^2,\allowbreak \ S_2^2=\dfrac{\displaystyle 1}{\displaystyle n}\displaystyle \sum_{i=1}^n(X_i-\bar X)^2,\allowbreak \ S_3^2=\dfrac{\displaystyle 1}{\displaystyle n-1}\displaystyle \sum_{i=1}^n(X_i-\mu)^2,\allowbreak \ S_4^2=\dfrac{\displaystyle 1}{\displaystyle n}\displaystyle \sum_{i=1}^n(X_i-\mu)^2\)，


\subsection*{试卷 IV（续）}

二、（5）（续）则服从自由度为 \(\displaystyle n-1\) 的 \(\displaystyle t\) 分布的随机变量是（　）。


（A）\(\displaystyle t=\dfrac{\displaystyle \bar X-\mu}{\displaystyle S_1/\sqrt{n-1}}\)。 （B）\(\displaystyle t=\dfrac{\displaystyle \bar X-\mu}{\displaystyle S_2/\sqrt{n-1}}\)。 （C）\(\displaystyle t=\dfrac{\displaystyle \bar X-\mu}{\displaystyle S_3/\sqrt n}\)。 （D）\(\displaystyle t=\dfrac{\displaystyle \bar X-\mu}{\displaystyle S_4/\sqrt n}\)。


三、（本题满分 6 分）计算二重积分 \(\displaystyle \iint_D(x+y)\,dx\,dy\)，其中 \(\displaystyle D=\{(x,\allowbreak y)\mid x^2+y^2\leq x+y+1\}\)。


四、（本题满分 5 分）设函数 \(\displaystyle y=y(x)\) 满足条件 \(\displaystyle y^{\prime\prime}+4y^{\prime}+4y=0,\allowbreak \ y(0)=2,\allowbreak \ y^{\prime}(0)=-4\)，求广义积分 \(\displaystyle \int_0^{+\infty}y(x)\,dx\)。


五、（本题满分 5 分）已知 \(\displaystyle f(x,\allowbreak y)=x^2\arctan\dfrac{\displaystyle y}{\displaystyle x}-y^2\arctan\dfrac{\displaystyle x}{\displaystyle y}\)，求 \(\displaystyle \dfrac{\displaystyle \partial^2 f}{\displaystyle \partial x\partial y}\)。


六、（本题满分 5 分）设函数 \(\displaystyle f(x)\) 可导，且 \(\displaystyle f(0)=0,\allowbreak  F(x)=\displaystyle \int_0^x t^{n-1}f(x^n-t^n)\,dt\)，求 \(\displaystyle \lim_{x\to0}\dfrac{\displaystyle F(x)}{\displaystyle x^{2n}}\)。


七、（本题满分 8 分）已知曲线 \(\displaystyle y=a\sqrt x\)（\(\displaystyle a>0\)）与曲线 \(\displaystyle y=\ln\sqrt x\) 在点 \(\displaystyle (x_0,\allowbreak y_0)\) 处有公共切线。求：（1）常数 \(\displaystyle a\) 及切点 \(\displaystyle (x_0,\allowbreak y_0)\)；（2）两曲线与 \(\displaystyle x\) 轴围成的平面图形绕 \(\displaystyle x\) 轴旋转所得旋转体的体积 \(\displaystyle V_x\)。


八、（本题满分 6 分）假设 \(\displaystyle f(x)\) 在 \(\displaystyle [a,\allowbreak +\infty)\) 上连续，\(\displaystyle f^{\prime\prime}(x)\) 在 \(\displaystyle (a,\allowbreak +\infty)\) 内存在且大于零，记 \(\displaystyle F(x)=\dfrac{\displaystyle f(x)-f(a)}{\displaystyle x-a}\)（\(\displaystyle x>a\)）。证明：\(\displaystyle F(x)\) 在 \(\displaystyle (a,\allowbreak +\infty)\) 内单调增加。


九、（本题满分 11 分）设线性方程组 \(\displaystyle \begin{cases}x_1+a_1x_2+a_1^2x_3=a_1^3,\allowbreak \\x_1+a_2x_2+a_2^2x_3=a_2^3,\allowbreak \\x_1+a_3x_2+a_3^2x_3=a_3^3,\allowbreak \\x_1+a_4x_2+a_4^2x_3=a_4^3.\end{cases}\)


（1）证明：若 \(\displaystyle a_1,\allowbreak a_2,\allowbreak a_3,\allowbreak a_4\) 两两不相等，则此线性方程组无解；


\subsection*{试卷 IV（续）}

九、（2）设 \(\displaystyle a_1=a_3=k,\allowbreak \ a_2=a_4=-k\)（\(\displaystyle k\ne0\)），且已知 \(\displaystyle \beta_1=\begin{pmatrix}-1\\1\\1\end{pmatrix},\allowbreak \ \beta_2=\begin{pmatrix}1\\1\\-1\end{pmatrix}\) 是该方程组的两个解，写出此方程组的通解。


十、（本题满分 8 分）设 \(\displaystyle A=\begin{pmatrix}0&0&1\\x&1&y\\1&0&0\end{pmatrix}\) 有三个线性无关的特征向量，求 \(\displaystyle x\) 和 \(\displaystyle y\) 应满足的条件。


十一、（本题满分 8 分）假设随机变量 \(\displaystyle X_1,\allowbreak X_2,\allowbreak X_3,\allowbreak X_4\) 相互独立，且同分布，\(\displaystyle P\{X_i=0\}=0.6,\allowbreak \ P\{X_i=1\}=0.4\)（\(\displaystyle i=1,\allowbreak 2,\allowbreak 3,\allowbreak 4\)）。求行列式 \(\displaystyle X=\begin{vmatrix}X_1&X_2\\X_3&X_4\end{vmatrix}\) 的概率分布。


十二、（本题满分 8 分）假设由自动线加工的某种零件的内径 \(\displaystyle X\)（毫米）服从正态分布 \(\displaystyle N(\mu,\allowbreak 1)\)，内径小于 10 或大于 12 为不合格品，其余为合格品。销售每件合格品获利，销售每件不合格品亏损。已知销售利润 \(\displaystyle T\)（单位：元）与销售零件的内径 \(\displaystyle X\) 有如下关系：\(\displaystyle T=\begin{cases}-1,\allowbreak &X<10,\allowbreak \\20,\allowbreak &10\leq X\leq12,\allowbreak \\-5,\allowbreak &X>12.\end{cases}\) 问平均内径 \(\displaystyle \mu\) 取何值时，销售一个零件的平均利润最大？


\subsection*{1994 年全国硕士研究生招生考试试题（试卷 V）}

\subsection*{一、填空题（本题共 5 小题，每小题 3 分，满分 15 分）}

（1）【同试卷 IV 第一、（1）题】


（2）【同试卷 IV 第一、（2）题】


（3）【同试卷 IV 第一、（3）题】


（4）【同试卷 IV 第一、（4）题】


（5）假设一批产品中一、二、三等品各占 60\%、30\%、10\%，从中随意取出一件，结果不是三等品，则取到的是一等品的概率为\_\_\_\_\_\_。


\subsection*{二、选择题（本题共 5 小题，每小题 3 分，满分 15 分）}

（1）【同试卷 IV 第二、（1）题】


（2）设函数 \(\displaystyle f(x)\) 在闭区间 \(\displaystyle [a,\allowbreak b]\) 上连续，且 \(\displaystyle f(x)>0\)，则方程 \(\displaystyle \int_a^x f(t)\,dt+\displaystyle \int_b^x\dfrac{\displaystyle 1}{\displaystyle f(t)}\,dt=0\) 在开区间 \(\displaystyle (a,\allowbreak b)\) 内的根有（　）。


（A）0 个。 （B）1 个。 （C）2 个。 （D）无穷多个。


（3）设 \(\displaystyle A,\allowbreak B\) 都是 \(\displaystyle n\) 阶非零矩阵，且 \(\displaystyle AB=O\)，则 \(\displaystyle A\) 和 \(\displaystyle B\) 的秩（　）。


（A）必有一个等于零。 （B）都小于 \(\displaystyle n\)。


\subsection*{试卷 V（续）}

二、（3）（续）（C）一个小于 \(\displaystyle n\)，一个等于 \(\displaystyle n\)。 （D）都等于 \(\displaystyle n\)。


（4）设有向量组 \(\displaystyle \alpha_1=(1,\allowbreak -1,\allowbreak 2,\allowbreak 4),\allowbreak \ \alpha_2=(0,\allowbreak 3,\allowbreak 1,\allowbreak 2),\allowbreak \ \alpha_3=(3,\allowbreak 0,\allowbreak 7,\allowbreak 14),\allowbreak \ \alpha_4=(1,\allowbreak -2,\allowbreak 2,\allowbreak 0),\allowbreak \ \alpha_5=(2,\allowbreak 1,\allowbreak 5,\allowbreak 10)\)，则该向量组的极大线性无关组是（　）。


（A）\(\displaystyle \alpha_1,\allowbreak \alpha_2,\allowbreak \alpha_3\)。 （B）\(\displaystyle \alpha_1,\allowbreak \alpha_2,\allowbreak \alpha_4\)。 （C）\(\displaystyle \alpha_1,\allowbreak \alpha_2,\allowbreak \alpha_5\)。 （D）\(\displaystyle \alpha_1,\allowbreak \alpha_2,\allowbreak \alpha_4,\allowbreak \alpha_5\)。


（5）【同试卷 IV 第二、（4）题】


三、（本题满分 5 分）求极限 \(\displaystyle \lim_{x\to\infty}\left[x-x^2\ln\left(1+\dfrac{\displaystyle 1}{\displaystyle x}\right)\right]\)。


四、（本题满分 5 分）【同试卷 IV 第五题】


五、（本题满分 6 分）已知 \(\displaystyle \dfrac{\displaystyle \sin x}{\displaystyle x}\) 是 \(\displaystyle f(x)\) 的一个原函数，求 \(\displaystyle \int x^3f^{\prime}(x)\,dx\)。


六、（本题满分 8 分）某养殖场饲养两种鱼，若甲种鱼放养 \(\displaystyle x\)（万尾），乙种鱼放养 \(\displaystyle y\)（万尾），收获时两种鱼的收获量分别为 \(\displaystyle (3-\alpha x-\beta y)x\) 和 \(\displaystyle (4-\beta x-2\alpha y)y\)（\(\displaystyle \alpha>\beta>0\)），求使产鱼总量最大的放养数。


七、（本题满分 8 分）已知曲线 \(\displaystyle y=a\sqrt x\)（\(\displaystyle a>0\)）与曲线 \(\displaystyle y=\ln\sqrt x\) 在点 \(\displaystyle (x_0,\allowbreak y_0)\) 处有公共切线，求：（1）常数 \(\displaystyle a\) 及切点 \(\displaystyle (x_0,\allowbreak y_0)\)；（2）两曲线与 \(\displaystyle x\) 轴围成的平面图形的面积 \(\displaystyle S\)。


八、（本题满分 7 分）【同试卷 IV 第六题】


九、（本题满分 8 分）设 \(\displaystyle \alpha_1,\allowbreak \alpha_2,\allowbreak \alpha_3\) 是齐次线性方程组 \(\displaystyle Ax=0\) 的一个基础解系。证明 \(\displaystyle \alpha_1+\alpha_2,\allowbreak \alpha_2+\alpha_3,\allowbreak \alpha_3+\alpha_1\) 也是该方程组的一个基础解系。


十、（本题满分 8 分）【同试卷 IV 第十题】


十一、（本题满分 7 分）假设随机变量 \(\displaystyle X\) 的概率密度为 \(\displaystyle f(x)=\begin{cases}2x,\allowbreak &0<x<1,\allowbreak \\0,\allowbreak &\text{其他}.\end{cases}\) 现在对 \(\displaystyle X\) 进行 \(\displaystyle n\) 次独立重复观测，以 \(\displaystyle V_n\) 表示观测值不大于 0.1 的次数，试求随机变量 \(\displaystyle V_n\) 的概率分布。


十二、（本题满分 8 分）【同试卷 IV 第十二题】

\end{document}
